%============================================================
%  Bd. 36 - Metrische Räume und Vollständigkeit
%============================================================

\documentclass[main.tex]{subfiles}

\ifSubfilesClassLoaded{
  \BandSetupStandalone{B36}
}{
}

\title{Bd. 36 - Metrische Räume und Vollständigkeit}
\author{Martin Kunze}
\date{}
\setcounter{file}{36}

\hypersetup{
  pdftitle={Bd. 36 - Metrische Räume und Vollständigkeit},
  pdfauthor={Martin Kunze}
}

\begin{document}

\title{Bd. 36 - Metrische Räume und Vollständigkeit}
\author{Martin Kunze}
\date{}
\hypersetup{pageanchor=false}
\maketitle
\hypersetup{pageanchor=true}
\setcounter{tocdepth}{3}
\tableofcontents
\setcounter{file}{36}
\renewcommand{\FormulaBandID}{36}

\chapter{Einleitung}

Band~19 hat auf den reellen Zahlen den Abstand
\(\RealDist{x}{y}=|x-y|\) konstruiert und seine vier Grundgesetze bewiesen.
Band~20 stellt die dabei benötigten Aussagen über endliche Anfangsabschnitte
bereit. Band~21 hat mit dem reellen Abstand Konvergenz und die
Cauchy-Eigenschaft reeller Folgen untersucht. In vielen dieser Beweise wurden
jedoch weder die Ordnung noch die Rechenoperationen der reellen Zahlen
benutzt. Benötigt wurden nur Positivität, Definitheit, Symmetrie und
Dreiecksungleichung des Abstands.

Dieser Band erhebt genau diese vier Eigenschaften zu Axiomen einer
\textbf{Metrik}. Dadurch werden Kugeln, Konvergenz, Cauchy-Folgen,
Vollständigkeit, Stetigkeit und Folgenkompaktheit auf beliebigen Mengen
verfügbar. Die bereits entwickelte reelle Theorie bleibt dabei das leitende
Beispiel und wird nicht ersetzt: Sie erscheint nun als Spezialfall des
metrischen Raums \((\mathbb R,d_{\mathbb R})\).

Eine Metrik setzt auf ihrer Trägermenge keine Addition und keine
Skalarmultiplikation voraus. Normen werden deshalb noch nicht als
Grundbegriff eingeführt. Sie gehören in die spätere Theorie der Vektorräume;
das Schlusskapitel ordnet bereits ein, wie jede Norm eine Metrik erzeugt.

Wie in Band~21 werden für reelle Zahlen die gewöhnlichen Symbole
\(0,1,+,-,\cdot,|\,\cdot\,|,<,\leq\) verwendet. Dies sind weiterhin nur die
dort vereinbarten Abkürzungen für die in Band~19 konstruierten Operationen.

\chapter{Metriken und offene Kugeln}

\section{Der metrische Grundbegriff}

\FormulaDefDeltaK[Metrik und metrischer Raum]{%
  \begin{aligned}[t]
  &d\colon X\times X\to\mathbb R\vdash
    \operatorname{Metrik}_{X}(d)\coloneqq{}\\[-2pt]
  &\quad\forall x,y,z\in X\,
    \Bigl(0\leq d(x,y)\land{}\\[-2pt]
  &\hspace{38mm}
      \bigl(d(x,y)=0\leftrightarrow x=y\bigr)\land{}\\[-2pt]
  &\hspace{38mm}
      d(x,y)=d(y,x)\land{}\\[-2pt]
  &\hspace{38mm}
      d(x,z)\leq d(x,y)+d(y,z)\Bigr).
  \end{aligned}%
}{MetricSpaceDef}{%
  \DeltaRow{Mengen}{X\dsep x\dsep y\dsep z}
  \DeltaRow{Funktionen}{d}
  \DeltaRow{Neues Prädikat}{\operatorname{Metrik}_{X}}
}

Gilt \(\operatorname{Metrik}_{X}(d)\), so heißt \(d\) eine
\textbf{Metrik auf \(X\)} und das Paar \((X,d)\) ein
\textbf{metrischer Raum}. Die vier Zeilen heißen Positivität,
Definitheit, Symmetrie und Dreiecksungleichung. Nur die Werte von \(d\) sind
reelle Zahlen; die Punkte von \(X\) brauchen keinerlei algebraische Struktur.

Die Definitheit trennt verschiedene Punkte: Der Abstand null ist nicht bloß
klein, sondern kennzeichnet Gleichheit. Die Dreiecksungleichung besagt, dass
ein Umweg über einen dritten Punkt den direkten Abstand nicht unterschreiten
kann.

\FormulaThmDeltaK[Umgekehrte Dreiecksungleichung]{%
  \operatorname{Metrik}_{X}(d)\dsep x,y,z\in X
  \vdash
  |d(x,z)-d(y,z)|\leq d(x,y)%
}{MetricReverseTriangleInequality}{%
  \DeltaRow{Mengen}{X\dsep x\dsep y\dsep z}
  \DeltaRow{Funktionen}{d}
}
\begin{proof}
Aus der Dreiecksungleichung folgen
\[
 d(x,z)\leq d(x,y)+d(y,z)
 \quad\text{und}\quad
 d(y,z)\leq d(y,x)+d(x,z)=d(x,y)+d(x,z).
\]
Nach Umstellen erhält man
\[
 d(x,z)-d(y,z)\leq d(x,y),
 \qquad
 d(y,z)-d(x,z)\leq d(x,y).
\]
Diese beiden Ungleichungen sind zusammen genau die behauptete
Abschätzung des Absolutbetrags.
\end{proof}

\section{Das reelle und das diskrete Beispiel}

Die in Band~19 eingeführte zweistellige Abstandsvorschrift besitzt für jedes
reelle Zahlenpaar genau einen reellen Wert. Das Prinzip der typisierten
Funktionsdefinition fasst diese Vorschrift daher zu einer Abbildung zusammen.

\FormulaDefDeltaK[Die reelle Abstandsfunktion]{%
  \begin{aligned}[t]
  &d_{\mathbb R}\colon\mathbb R\times\mathbb R\to\mathbb R,\\[-2pt]
  &x,y\in\mathbb R\vdash
    d_{\mathbb R}(x,y)\coloneqq|x-y|.
  \end{aligned}%
}{RealMetricDistanceFunctionDef}{%
  \DeltaRow{Mengen}{\mathbb R\dsep x\dsep y}
  \DeltaRow{Neue Funktion}{d_{\mathbb R}}
}

\FormulaThmDeltaK[Der reelle Abstand ist eine Metrik]{%
  \operatorname{Metrik}_{\mathbb R}(d_{\mathbb R})%
}{RealDistanceIsMetric}{%
  \DeltaRow{Mengen}{\mathbb R}
  \DeltaRow{Funktionen}{d_{\mathbb R}}
}
\begin{proof}
Die Funktionstypisierung stammt aus
\FormulaRefAuto{RealMetricDistanceFunctionDef}[def].

Nach
\FormulaRefAuto{RealDistance}[def] ist
\(d_{\mathbb R}(x,y)=|x-y|\). Positivität, Definitheit, Symmetrie und
Dreiecksungleichung sind genau die vier Aussagen aus
\FormulaRefAuto{RealDistanceLaws}[thm].
\end{proof}

Das nächste Beispiel zeigt, dass eine Metrik weder aus Zahlen noch aus einer
geometrischen Konstruktion stammen muss.

\FormulaDefDeltaK[Diskrete Metrik]{%
  \begin{aligned}[t]
  &\delta_X\colon X\times X\to\mathbb R,\\[-2pt]
  &x,y\in X\vdash\delta_X(x,y)\coloneqq
  \begin{cases}
    0,&x=y,\\
    1,&x\neq y,
  \end{cases}
  \end{aligned}%
}{DiscreteMetricDef}{%
  \DeltaRow{Mengen}{X\dsep x\dsep y}
  \DeltaRow{Neue Funktion}{\delta_X}
}

\FormulaThmDeltaK[Die diskrete Metrik ist eine Metrik]{%
  \operatorname{Metrik}_{X}(\delta_X)%
}{DiscreteMetricIsMetric}{%
  \DeltaRow{Mengen}{X}
  \DeltaRow{Funktionen}{\delta_X}
}
\begin{proof}
Die Werte \(0\) und \(1\) sind nichtnegativ. Nach Definition gilt
\(\delta_X(x,y)=0\) genau für \(x=y\), und die Fallunterscheidung ist
symmetrisch in \(x\) und \(y\). Für die Dreiecksungleichung ist nur der Fall
\(x\neq z\) zu betrachten. Dann ist die linke Seite gleich \(1\). Wären
beide Summanden rechts null, so wäre \(x=y=z\), ein Widerspruch. Daher ist
mindestens einer der beiden Summanden gleich \(1\), und die Ungleichung gilt.
\end{proof}

\section{Teilmengen und eingeschränkte Metriken}

\FormulaDefDeltaK[Teilraummetrik]{%
  \operatorname{Metrik}_{X}(d)\dsep Y\subseteq X
  \vdash
  d_Y\coloneqq d\restriction_{Y\times Y}%
}{MetricSubspaceRestrictionDef}{%
  \DeltaRow{Mengen}{X\dsep Y}
  \DeltaRow{Funktionen}{d\dsep d_Y}
}

\FormulaThmDeltaK[Einschränkung einer Metrik]{%
  \operatorname{Metrik}_{X}(d)\dsep Y\subseteq X
  \vdash
  \operatorname{Metrik}_{Y}(d_Y)%
}{MetricSubspaceRestrictionIsMetric}{%
  \DeltaRow{Mengen}{X\dsep Y}
  \DeltaRow{Funktionen}{d\dsep d_Y}
}
\begin{proof}
Für Punkte aus \(Y\) stimmen alle Werte von \(d_Y\) mit den entsprechenden
Werten von \(d\) überein. Daher vererben sich alle vier Metrikaxiome an die
Einschränkung.
\end{proof}

Wenn von einer Teilmenge eines metrischen Raums als metrischem Raum gesprochen
wird, ist ohne weiteren Zusatz stets diese Teilraummetrik gemeint.

\section{Offene Kugeln und offene Mengen}

\FormulaDefDeltaK[Offene metrische Kugel]{%
  \begin{aligned}[t]
  &\operatorname{Metrik}_{X}(d)\dsep x\in X\dsep r\in\mathbb R\dsep0<r
  \vdash\\[-2pt]
  &\quad\MetricBall{d}{x}{r}\coloneqq
    \{y\in X\mid d(y,x)<r\}.
  \end{aligned}%
}{MetricOpenBallDef}{%
  \DeltaRow{Mengen}{X\dsep x\dsep y\dsep r}
  \DeltaRow{Funktionen}{d}
  \DeltaRow{Neue Schreibweise}{\MetricBall{d}{x}{r}}
}

Der Radius ist positiv und der Mittelpunkt gehört wegen
\(d(x,x)=0<r\) zu seiner Kugel. In \(\mathbb R\) sind diese Kugeln genau die
bereits bekannten offenen Intervalle.

\FormulaThmDeltaK[Reelle Kugeln sind Epsilon-Umgebungen]{%
  x,r\in\mathbb R\dsep0<r
  \vdash
  \MetricBall{d_{\mathbb R}}{x}{r}=B_r(x)%
}{RealMetricBallIsEpsilonNeighborhood}{%
  \DeltaRow{Mengen}{x\dsep r}
  \DeltaRow{Funktionen}{d_{\mathbb R}}
}
\begin{proof}
Beide Mengen bestehen nach
\FormulaRefAuto{MetricOpenBallDef}[def] und
\FormulaRefAuto{RealEpsilonNeighborhood}[def] genau aus den reellen Zahlen
\(y\) mit \(d_{\mathbb R}(y,x)<r\).
\end{proof}

\FormulaThmDeltaK[Kugelverkleinerung]{%
  \begin{aligned}[t]
  &\operatorname{Metrik}_{X}(d)\dsep x\in X\dsep r>0\dsep
    y\in\MetricBall{d}{x}{r}\vdash\\[-2pt]
  &\quad \Bigl(r-d(x,y)>0\land
    \MetricBall{d}{y}{r-d(x,y)}
      \subseteq\MetricBall{d}{x}{r}\Bigr).
  \end{aligned}%
}{MetricBallShrinking}{%
  \DeltaRow{Mengen}{X\dsep x\dsep y\dsep r}
  \DeltaRow{Funktionen}{d}
}
\begin{proof}
Aus \(y\in B_r^d(x)\) folgt \(d(x,y)<r\), also
\(r-d(x,y)>0\). Ist
\(z\in B_{r-d(x,y)}^d(y)\), so liefert die Dreiecksungleichung
\[
 d(x,z)\leq d(x,y)+d(y,z)
 <d(x,y)+r-d(x,y)=r.
\]
Somit liegt \(z\) in \(B_r^d(x)\).
\end{proof}

\FormulaDefDeltaK[Offene und abgeschlossene Teilmenge]{%
  \begin{aligned}[t]
  &\operatorname{Metrik}_{X}(d)\dsep U,A\subseteq X\vdash\\[-2pt]
  &\quad\operatorname{Offen}_{d}(U)\coloneqq
    \forall x\in U\,\exists r\in\mathbb R\,
      \bigl(0<r\land\MetricBall{d}{x}{r}\subseteq U\bigr),\\[-2pt]
  &\quad\operatorname{Abgeschlossen}_{d}(A)\coloneqq
    \operatorname{Offen}_{d}(X\setminus A).
  \end{aligned}%
}{MetricOpenClosedSetDef}{%
  \DeltaRow{Mengen}{X\dsep U\dsep A\dsep x\dsep r}
  \DeltaRow{Funktionen}{d}
  \DeltaRow{Neue Prädikate}{\operatorname{Offen}_{d}\dsep
    \operatorname{Abgeschlossen}_{d}}
}

\FormulaThmDeltaK[Offene Kugeln sind offen]{%
  \operatorname{Metrik}_{X}(d)\dsep x\in X\dsep r>0
  \vdash
  \operatorname{Offen}_{d}\!\left(\MetricBall{d}{x}{r}\right)%
}{MetricOpenBallsAreOpen}{%
  \DeltaRow{Mengen}{X\dsep x\dsep r}
  \DeltaRow{Funktionen}{d}
}
\begin{proof}
Für jeden Punkt \(y\in B_r^d(x)\) liefert
\FormulaRefAuto{MetricBallShrinking}[thm] den positiven Radius
\(r-d(x,y)\) einer um \(y\) liegenden Kugel, die ganz in
\(B_r^d(x)\) enthalten ist. Dies ist genau die Offenheitsbedingung.
\end{proof}

\chapter{Folgen in metrischen Räumen}

\section{Metrische Konvergenz}

\FormulaDefDeltaK[Konvergenz in einem metrischen Raum]{%
  \begin{aligned}[t]
  &\operatorname{Metrik}_{X}(d)\dsep a\colon\mathbb N\to X\dsep x\in X
  \vdash\\[-2pt]
  &\quad\bigl(a_n\MetricTo{d}x\bigr)\coloneqq
    \forall\varepsilon\in\mathbb R\,
    \Bigl(0<\varepsilon\rightarrow
      \exists N\in\mathbb N\,\forall n\in\mathbb N\\[-2pt]
  &\hspace{63mm}
      (N\leq n\rightarrow d(a_n,x)<\varepsilon)\Bigr).
  \end{aligned}%
}{MetricSequenceConvergenceDef}{%
  \DeltaRow{Mengen}{X\dsep x\dsep\varepsilon\dsep N\dsep n}
  \DeltaRow{Funktionen}{d\dsep a}
  \DeltaRow{Neue Schreibweise}{a_n\MetricTo{d}x}
}

Die Aussage bedeutet: Jede noch so kleine Kugel um \(x\) enthält alle bis
auf endlich viele Folgenglieder. Die Folge darf dieselben Werte mehrfach
annehmen; entscheidend ist ihre Reihenfolge, nicht nur ihre Wertemenge.

\FormulaThmDeltaK[Übereinstimmung mit der reellen Konvergenz]{%
  a\colon\mathbb N\to\mathbb R\dsep L\in\mathbb R
  \vdash
  \bigl(a_n\MetricTo{d_{\mathbb R}}L\bigr)
  \leftrightarrow
  \bigl(a_n\longrightarrow L\bigr)%
}{RealAndMetricSequenceConvergenceAgree}{%
  \DeltaRow{Mengen}{L}
  \DeltaRow{Funktionen}{a\dsep d_{\mathbb R}}
}
\begin{proof}
Nach \FormulaRefAuto{RealDistance}[def] gilt
\(d_{\mathbb R}(a_n,L)=|a_n-L|\). Nach Einsetzen dieser Gleichheit stimmen
\FormulaRefAuto{MetricSequenceConvergenceDef}[def] und
\FormulaRefAuto{RealSequenceConvergenceDef}[def] Wort für Wort überein.
\end{proof}

\FormulaThmDeltaK[Konvergenz in einem metrischen Teilraum]{%
  \begin{aligned}[t]
  &\operatorname{Metrik}_{X}(d)\dsep Y\subseteq X\dsep
    a\colon\mathbb N\to Y\dsep y\in Y\vdash\\[-2pt]
  &\quad
    \bigl(a_n\MetricTo{d_Y}y\bigr)
    \leftrightarrow
    \bigl(a_n\MetricTo{d}y\bigr).
  \end{aligned}%
}{MetricSubspaceConvergenceAgree}{%
  \DeltaRow{Mengen}{X\dsep Y\dsep y\dsep\varepsilon\dsep N\dsep n}
  \DeltaRow{Funktionen}{d\dsep d_Y\dsep a}
}
\begin{proof}
Da \(Y\subseteq X\), kann \(a\) zugleich als Folge mit Werten in \(X\)
aufgefasst werden. Für alle Folgenglieder gilt nach
\FormulaRefAuto{MetricSubspaceRestrictionDef}[def]
\(d_Y(a_n,y)=d(a_n,y)\). Nach Einsetzen dieser Gleichheit sind beide
Konvergenzbedingungen identisch.
\end{proof}

\FormulaThmDeltaK[Eindeutigkeit metrischer Grenzwerte]{%
  \operatorname{Metrik}_{X}(d)\dsep a\colon\mathbb N\to X\dsep x,y\in X
  \dsep a_n\MetricTo{d}x\dsep a_n\MetricTo{d}y
  \vdash x=y%
}{MetricSequenceLimitUnique}{%
  \DeltaRow{Mengen}{X\dsep x\dsep y\dsep\varepsilon\dsep N\dsep n}
  \DeltaRow{Funktionen}{d\dsep a}
}
\begin{proof}
Angenommen, \(x\neq y\). Dann ist
\(\varepsilon=d(x,y)/3>0\). Für einen hinreichend großen gemeinsamen
Index \(n\) gelten
\(d(a_n,x)<\varepsilon\) und \(d(a_n,y)<\varepsilon\). Mit Symmetrie und
Dreiecksungleichung folgt
\[
 d(x,y)\leq d(x,a_n)+d(a_n,y)
 <2\varepsilon=\frac23d(x,y),
\]
ein Widerspruch.
\end{proof}

\FormulaThmDeltaK[Konstante Folgen konvergieren]{%
  \operatorname{Metrik}_{X}(d)\dsep x\in X
  \vdash
  (\operatorname{const}_{\mathbb N,x})_n\MetricTo{d}x%
}{ConstantMetricSequenceConverges}{%
  \DeltaRow{Mengen}{X\dsep x\dsep n}
  \DeltaRow{Funktionen}{d\dsep\operatorname{const}_{\mathbb N,x}}
}
\begin{proof}
Jedes Folgenglied ist \(x\), und \(d(x,x)=0<\varepsilon\) für jeden
positiven Radius \(\varepsilon\).
\end{proof}

\FormulaThmDeltaK[Teilfolgen erhalten metrische Grenzwerte]{%
  \begin{aligned}[t]
  &\operatorname{Metrik}_{X}(d)\dsep a\colon\mathbb N\to X\dsep x\in X
    \dsep a_n\MetricTo{d}x\dsep
    \varphi\colon\mathbb N\to\mathbb N\dsep{}\\[-2pt]
  &\quad\forall n\in\mathbb N\,
      \bigl(\varphi(n)<\varphi(\Succ(n))\bigr)
    \vdash a_{\varphi(n)}\MetricTo{d}x.
  \end{aligned}%
}{MetricSubsequencePreservesLimit}{%
  \DeltaRow{Mengen}{X\dsep x\dsep n}
  \DeltaRow{Funktionen}{d\dsep a\dsep\varphi}
}
\begin{proof}
Eine streng wachsende Folge natürlicher Indizes erfüllt
\(n\leq\varphi(n)\). Sobald \(n\geq N\) gilt, ist daher auch
\(\varphi(n)\geq N\), und die Abschätzung der Ausgangsfolge ist anwendbar.
\end{proof}

\FormulaThmDeltaK[Endliche Änderungen erhalten metrische Grenzwerte]{%
  \begin{aligned}[t]
  &\operatorname{Metrik}_{X}(d)\dsep a,b\colon\mathbb N\to X\dsep
    \exists N_0\in\mathbb N\,\forall n\in\mathbb N\,
      (N_0\leq n\rightarrow a_n=b_n)\vdash\\[-2pt]
  &\quad\forall x\in X\,
    \bigl(a_n\MetricTo{d}x\leftrightarrow b_n\MetricTo{d}x\bigr).
  \end{aligned}%
}{MetricFiniteModificationPreservesLimit}{%
  \DeltaRow{Mengen}{X\dsep x\dsep N_0\dsep N\dsep n}
  \DeltaRow{Funktionen}{d\dsep a\dsep b}
}
\begin{proof}
Zu einem Konvergenzindex \(N\) der einen Folge verwendet man für die andere
den größeren der beiden Indizes \(N\) und \(N_0\). Ab diesem Index stimmen
die Folgenglieder und damit ihre Abstände von \(x\) überein.
\end{proof}

\FormulaThmDeltaK[Restfolgen haben denselben metrischen Grenzwert]{%
  \operatorname{Metrik}_{X}(d)\dsep a\colon\mathbb N\to X\dsep
  r\in\mathbb N\dsep x\in X\vdash
  a_n\MetricTo{d}x
  \leftrightarrow
  a^{\langle r\rangle}_n\MetricTo{d}x%
}{MetricSequenceTailSameLimit}{%
  \DeltaRow{Mengen}{X\dsep x\dsep r\dsep n}
  \DeltaRow{Funktionen}{d\dsep a\dsep a^{\langle r\rangle}}
}
\begin{proof}
Konvergiert \(a\), so genügt für die Restfolge derselbe Index, denn
\(n+r\geq n\). Konvergiert umgekehrt die Restfolge ab dem Index \(N\), so
gilt die benötigte Abschätzung für \(a_n\) ab \(N+r\): Dann ist
\(n=k+r\) mit \(k\geq N\), und
\(a_n=a_{k+r}=a^{\langle r\rangle}_k\).
\end{proof}

\section{Metrische Beschränktheit}

\FormulaDefDeltaK[Beschränkte Teilmenge und beschränkte Folge]{%
  \begin{aligned}[t]
  &\operatorname{Metrik}_{X}(d)\dsep A\subseteq X\vdash\\[-2pt]
  &\quad\operatorname{Beschr}_{d}(A)\coloneqq
    \Bigl(A=\varnothing\lor
      \exists c\in X\,\exists r\in\mathbb R\,
      \bigl(0<r\land A\subseteq\MetricBall{d}{c}{r}\bigr)\Bigr),\\[-2pt]
  &\operatorname{Metrik}_{X}(d)\dsep a\colon\mathbb N\to X\vdash\\[-2pt]
  &\quad\operatorname{Beschr}_{d}(a)\coloneqq
    \operatorname{Beschr}_{d}\bigl(a[\mathbb N]\bigr).
  \end{aligned}%
}{MetricBoundedSetSequenceDef}{%
  \DeltaRow{Mengen}{X\dsep A\dsep c\dsep r}
  \DeltaRow{Funktionen}{d\dsep a}
  \DeltaRow{Neues Prädikat}{\operatorname{Beschr}_{d}}
}

Die ausdrückliche Leermengenalternative sorgt dafür, dass auch die leere
Teilmenge eines leeren metrischen Raums beschränkt ist. Für nichtleere Mengen
ist die Definition gleichbedeutend mit der üblichen Forderung, dass alle
Punkte in einer hinreichend großen Kugel liegen.

\FormulaThmDeltaK[Konvergente metrische Folgen sind beschränkt]{%
  \operatorname{Metrik}_{X}(d)\dsep a\colon\mathbb N\to X\dsep x\in X
  \dsep a_n\MetricTo{d}x
  \vdash\operatorname{Beschr}_{d}(a)%
}{ConvergentMetricSequenceBounded}{%
  \DeltaRow{Mengen}{X\dsep x\dsep N\dsep n\dsep r}
  \DeltaRow{Funktionen}{d\dsep a}
}
\begin{proof}
Für den Radius \(1\) gibt es einen Index \(N\), ab dem alle Glieder in
\(B_1^d(x)\) liegen. Die Menge der vorherigen Glieder ist endlich. Daher sind
auch die endlich vielen reellen Abstände \(d(a_n,x)\) mit \(n<N\) gemeinsam
nach oben beschränkt. Wählt man \(r>1\), das auch größer als alle diese
Abstände ist, so liegt die gesamte Wertemenge in \(B_r^d(x)\).
\end{proof}

\FormulaThmDeltaK[Reelle und metrische Beschränktheit stimmen überein]{%
  a\colon\mathbb N\to\mathbb R
  \vdash
  \operatorname{Beschr}_{d_{\mathbb R}}(a)
  \leftrightarrow
  \operatorname{Beschr}(a)%
}{RealAndMetricSequenceBoundednessAgree}{%
  \DeltaRow{Mengen}{\mathbb R}
  \DeltaRow{Funktionen}{a\dsep d_{\mathbb R}}
}
\begin{proof}
Da \(0\in\mathbb N\) und \(a_0\in a[\mathbb N]\), ist die Wertemenge der
Folge nichtleer. In der metrischen Beschränktheitsdefinition kann daher die
Leermengenalternative nicht eintreten.

Ist \(|a_n|\leq C\) für alle \(n\), so liegt die Wertemenge in
\(B_{C+1}^{d_{\mathbb R}}(0)\). Liegt sie umgekehrt in
\(B_r^{d_{\mathbb R}}(c)\), so gilt mit der Dreiecksungleichung
\[
 |a_n|=d_{\mathbb R}(a_n,0)
 \leq d_{\mathbb R}(a_n,c)+d_{\mathbb R}(c,0)
 <r+|c|.
\]
Somit besitzt die Folge auch eine Schranke im Sinne von
\FormulaRefAuto{BoundedRealSequenceDef}[def].
\end{proof}

\chapter{Cauchy-Folgen und Vollständigkeit}

\section{Cauchy-Folgen}

\FormulaDefDeltaK[Cauchy-Folge in einem metrischen Raum]{%
  \begin{aligned}[t]
  &\operatorname{Metrik}_{X}(d)\dsep a\colon\mathbb N\to X\vdash
    \operatorname{Cauchy}_{d}(a)\coloneqq{}\\[-2pt]
  &\quad\forall\varepsilon\in\mathbb R\,
    \Bigl(0<\varepsilon\rightarrow
      \exists N\in\mathbb N\,\forall m,n\in\mathbb N\\[-2pt]
  &\hspace{50mm}
      \bigl((N\leq m\land N\leq n)
      \rightarrow d(a_m,a_n)<\varepsilon\bigr)\Bigr).
  \end{aligned}%
}{MetricCauchySequenceDef}{%
  \DeltaRow{Mengen}{X\dsep\varepsilon\dsep N\dsep m\dsep n}
  \DeltaRow{Funktionen}{d\dsep a}
  \DeltaRow{Neues Prädikat}{\operatorname{Cauchy}_{d}}
}

Die Bedingung nennt keinen Grenzwert. Sie prüft nur, ob späte Glieder
untereinander beliebig nahe liegen. Genau deshalb kann sie auch in einem Raum
formuliert werden, in dem der erwartete Grenzpunkt fehlt.

\FormulaThmDeltaK[Übereinstimmung der beiden reellen Cauchy-Begriffe]{%
  a\colon\mathbb N\to\mathbb R
  \vdash
  \operatorname{Cauchy}_{d_{\mathbb R}}(a)
  \leftrightarrow
  \operatorname{Cauchy}(a)%
}{RealAndMetricCauchySequenceAgree}{%
  \DeltaRow{Mengen}{\mathbb R}
  \DeltaRow{Funktionen}{a\dsep d_{\mathbb R}}
}
\begin{proof}
Nach \FormulaRefAuto{RealDistance}[def] gilt
\(d_{\mathbb R}(a_m,a_n)=|a_m-a_n|\). Mit dieser Gleichheit sind
\FormulaRefAuto{MetricCauchySequenceDef}[def] und
\FormulaRefAuto{RealCauchySequenceDef}[def] dieselbe Bedingung.
\end{proof}

\FormulaThmDeltaK[Cauchy-Eigenschaft in einem metrischen Teilraum]{%
  \begin{aligned}[t]
  &\operatorname{Metrik}_{X}(d)\dsep Y\subseteq X\dsep
    a\colon\mathbb N\to Y\vdash\\[-2pt]
  &\quad
    \operatorname{Cauchy}_{d_Y}(a)
    \leftrightarrow
    \operatorname{Cauchy}_{d}(a).
  \end{aligned}%
}{MetricSubspaceCauchyAgree}{%
  \DeltaRow{Mengen}{X\dsep Y\dsep\varepsilon\dsep N\dsep m\dsep n}
  \DeltaRow{Funktionen}{d\dsep d_Y\dsep a}
}
\begin{proof}
Da \(Y\subseteq X\), kann \(a\) zugleich als Folge mit Werten in \(X\)
aufgefasst werden. Für alle \(m,n\in\mathbb N\) gilt
\(d_Y(a_m,a_n)=d(a_m,a_n)\). Daher stimmen die beiden Cauchy-Bedingungen
nach \FormulaRefAuto{MetricCauchySequenceDef}[def] vollständig überein.
\end{proof}

\FormulaThmDeltaK[Konvergenz impliziert die metrische Cauchy-Eigenschaft]{%
  \operatorname{Metrik}_{X}(d)\dsep a\colon\mathbb N\to X\dsep x\in X
  \dsep a_n\MetricTo{d}x
  \vdash\operatorname{Cauchy}_{d}(a)%
}{MetricConvergentSequenceIsCauchy}{%
  \DeltaRow{Mengen}{X\dsep x\dsep\varepsilon\dsep N\dsep m\dsep n}
  \DeltaRow{Funktionen}{d\dsep a}
}
\begin{proof}
Zu \(\varepsilon>0\) wähle \(N\) so, dass
\(d(a_k,x)<\varepsilon/2\) für alle \(k\geq N\). Für \(m,n\geq N\)
folgt
\[
 d(a_m,a_n)\leq d(a_m,x)+d(x,a_n)<\varepsilon.
\]
\end{proof}

\FormulaThmDeltaK[Metrische Cauchy-Folgen sind beschränkt]{%
  \operatorname{Metrik}_{X}(d)\dsep a\colon\mathbb N\to X\dsep
  \operatorname{Cauchy}_{d}(a)
  \vdash\operatorname{Beschr}_{d}(a)%
}{MetricCauchySequenceBounded}{%
  \DeltaRow{Mengen}{X\dsep N\dsep n\dsep r}
  \DeltaRow{Funktionen}{d\dsep a}
}
\begin{proof}
Für \(\varepsilon=1\) liegen ab einem Index \(N\) alle Glieder in
\(B_1^d(a_N)\). Die endlich vielen vorherigen Glieder besitzen gemeinsam mit
dieser Kugel eine hinreichend große Kugel um \(a_N\), genau wie im Beweis von
\FormulaRefAuto{ConvergentMetricSequenceBounded}[thm].
\end{proof}

\FormulaThmDeltaK[Cauchy-Folge mit konvergenter Teilfolge]{%
  \begin{aligned}[t]
  &\operatorname{Metrik}_{X}(d)\dsep a\colon\mathbb N\to X\dsep
    \operatorname{Cauchy}_{d}(a)\dsep x\in X\dsep
    \varphi\colon\mathbb N\to\mathbb N\dsep{}\\[-2pt]
  &\quad\forall n\in\mathbb N\,
      \bigl(\varphi(n)<\varphi(\Succ(n))\bigr)
    \dsep a_{\varphi(n)}\MetricTo{d}x
    \vdash a_n\MetricTo{d}x.
  \end{aligned}%
}{MetricCauchyWithConvergentSubsequence}{%
  \DeltaRow{Mengen}{X\dsep x\dsep\varepsilon\dsep N\dsep n}
  \DeltaRow{Funktionen}{d\dsep a\dsep\varphi}
}
\begin{proof}
Zu \(\varepsilon>0\) wähle \(N_1\) für die Cauchy-Bedingung mit
\(\varepsilon/2\) und \(N_2\) so, dass
\(d(a_{\varphi(k)},x)<\varepsilon/2\) für \(k\geq N_2\). Wähle ein
\(k\geq N_2\) mit \(\varphi(k)\geq N_1\). Für jedes \(n\geq N_1\) gilt
dann
\[
 d(a_n,x)\leq d(a_n,a_{\varphi(k)})
               +d(a_{\varphi(k)},x)<\varepsilon.
\]
\end{proof}

\section{Vollständige metrische Räume}

\FormulaDefDeltaK[Vollständiger metrischer Raum]{%
  \begin{aligned}[t]
  &\operatorname{Metrik}_{X}(d)\vdash
    \operatorname{Vollst}(X,d)\coloneqq{}\\[-2pt]
  &\quad\forall a\colon\mathbb N\to X\,
    \Bigl(\operatorname{Cauchy}_{d}(a)\rightarrow
      \exists x\in X\,(a_n\MetricTo{d}x)\Bigr).
  \end{aligned}%
}{CompleteMetricSpaceDef}{%
  \DeltaRow{Mengen}{X\dsep x}
  \DeltaRow{Funktionen}{d\dsep a}
  \DeltaRow{Neues Prädikat}{\operatorname{Vollst}}
}

Vollständigkeit bedeutet, dass im Raum kein Grenzpunkt fehlt, den seine
eigene Abstandsmessung durch eine Cauchy-Folge ankündigt. Sie ist von der
Ordnungsvollständigkeit zu unterscheiden: Die Vollständigkeit eines
geordneten Körpers handelt von Suprema, die metrische Vollständigkeit von
Cauchy-Folgen.

\FormulaThmDeltaK[Die reellen Zahlen sind metrisch vollständig]{%
  \operatorname{Vollst}(\mathbb R,d_{\mathbb R})%
}{RealMetricSpaceComplete}{%
  \DeltaRow{Mengen}{\mathbb R}
  \DeltaRow{Funktionen}{d_{\mathbb R}}
}
\begin{proof}
Nach \FormulaRefAuto{RealAndMetricSequenceConvergenceAgree}[thm] stimmt die
metrische Konvergenz mit der reellen Konvergenz aus Band~21 überein. Ebenso
stimmt nach
\FormulaRefAuto{RealAndMetricCauchySequenceAgree}[thm] die metrische
Cauchy-Bedingung mit der reellen überein. Das in
\FormulaRefAuto{RealSequenceCauchyCriterion}[thm] bewiesene reelle
Cauchy-Kriterium liefert daher für jede metrische Cauchy-Folge einen reellen
Grenzwert.

Das Cauchy-Kriterium aus Band~21 beruht seinerseits auf der Existenz von
Infima und Suprema aus
\FormulaRefAuto{RealInfimumCharacterization}[thm] und
\FormulaRefAuto{RealLeastUpperBoundProperty}[thm]. Damit ist die
Abhängigkeitskette ausdrücklich sichtbar:
Ordnungsvollständigkeit der reellen Zahlen liefert ihr Cauchy-Kriterium und
dieses wiederum die metrische Vollständigkeit von
\((\mathbb R,d_{\mathbb R})\).
\end{proof}

\FormulaThmDeltaK[Diskrete metrische Räume sind vollständig]{%
  \operatorname{Vollst}(X,\delta_X)%
}{DiscreteMetricSpaceComplete}{%
  \DeltaRow{Mengen}{X}
  \DeltaRow{Funktionen}{\delta_X}
}
\begin{proof}
Sei \(a\colon\mathbb N\to X\) eine Cauchy-Folge. Für
\(\varepsilon=1/2\) gibt es ein \(N\), so dass
\(\delta_X(a_m,a_n)<1/2\) für alle \(m,n\geq N\). Nach Definition der
diskreten Metrik ist dann \(a_m=a_n\). Die Folge ist also ab \(N\) konstant
gleich \(a_N\) und konvergiert gegen diesen Punkt.
\end{proof}

Für das folgende Gegenbeispiel benötigen wir eine elementare Nullfolge.

\FormulaThmDeltaK[Kehrwerte positiver natürlicher Zahlen bilden eine Nullfolge]{%
  \begin{aligned}[t]
  &r\colon\mathbb N\to\mathbb R\dsep
    \forall n\in\mathbb N\,
      \left(r_n=\frac{1}{n+1}\right)
    \vdash r_n\longrightarrow0.
  \end{aligned}%
}{ReciprocalNaturalSequenceConvergesToZero}{%
  \DeltaRow{Mengen}{n\dsep N}
  \DeltaRow{Funktionen}{r}
}
\begin{proof}
Sei \(\varepsilon>0\). Nach der archimedischen Eigenschaft der reellen
Zahlen aus \FormulaRefAuto{RealArchimedeanProperty}[thm] gibt es ein
\(N\in\mathbb N\) mit \(1/\varepsilon<N+1\). Für \(n\geq N\) gilt dann
\[
 0<\frac1{n+1}\leq\frac1{N+1}<\varepsilon.
\]
Damit ist die Konvergenzdefinition erfüllt.
\end{proof}

\FormulaThmDeltaK[Die punktierte reelle Gerade ist nicht vollständig]{%
  \neg\operatorname{Vollst}
    \bigl(\mathbb R\setminus\{0\},d_{\mathbb R\setminus\{0\}}\bigr)%
}{PuncturedRealLineIncomplete}{%
  \DeltaRow{Mengen}{\mathbb R\dsep\mathbb R\setminus\{0\}\dsep 0}
  \DeltaRow{Funktionen}{d_{\mathbb R}\dsep
    d_{\mathbb R\setminus\{0\}}\dsep a}
}
\begin{proof}
Setze \(A=\mathbb R\setminus\{0\}\) und
\(a_n=1/(n+1)\). Dann ist \(a\colon\mathbb N\to A\). Nach
\FormulaRefAuto{ReciprocalNaturalSequenceConvergesToZero}[thm] und
\FormulaRefAuto{RealAndMetricSequenceConvergenceAgree}[thm] gilt
\(a_n\MetricTo{d_{\mathbb R}}0\). Daher ist \(a\) nach
\FormulaRefAuto{MetricConvergentSequenceIsCauchy}[thm] eine Cauchy-Folge in
\((\mathbb R,d_{\mathbb R})\), nach
\FormulaRefAuto{MetricSubspaceCauchyAgree}[thm] also auch in
\((A,d_A)\).

Wäre dieser Teilraum vollständig, so gäbe es ein \(x\in A\) mit
\(a_n\MetricTo{d_A}x\). Nach
\FormulaRefAuto{MetricSubspaceConvergenceAgree}[thm] konvergierte \(a\) im
reellen metrischen Raum auch gegen \(x\). Die Eindeutigkeit metrischer
Grenzwerte ergäbe aus der gleichzeitigen Konvergenz gegen \(0\) die Gleichheit
\(x=0\), im Widerspruch zu \(x\in A\).
\end{proof}

Das Gegenbeispiel entfernt genau den Punkt, den die Cauchy-Folge als Grenzwert
ankündigt. Der folgende Satz zeigt, welche Teilmengen eines vollständigen
Raums dieses Problem vermeiden.

\FormulaThmDeltaK[Abgeschlossene Mengen enthalten ihre Folgengrenzwerte]{%
  \begin{aligned}[t]
  &\operatorname{Metrik}_{X}(d)\dsep A\subseteq X\dsep
    \operatorname{Abgeschlossen}_{d}(A)\dsep
    a\colon\mathbb N\to A\dsep x\in X\dsep{}\\[-2pt]
  &\quad a_n\MetricTo{d}x\vdash x\in A.
  \end{aligned}%
}{MetricClosedSetContainsSequenceLimits}{%
  \DeltaRow{Mengen}{X\dsep A\dsep x\dsep r\dsep N}
  \DeltaRow{Funktionen}{d\dsep a}
}
\begin{proof}
Wäre \(x\notin A\), so läge \(x\) in der offenen Menge
\(X\setminus A\). Dann gäbe es ein \(r>0\) mit
\(\MetricBall{d}{x}{r}\subseteq X\setminus A\). Aus der Konvergenz folgte
aber \(a_N\in\MetricBall{d}{x}{r}\) für einen hinreichend großen Index
\(N\), obwohl \(a_N\in A\) gilt. Dieser Widerspruch zeigt \(x\in A\).
\end{proof}

\FormulaThmDeltaK[Abgeschlossene Teilräume vollständiger Räume]{%
  \begin{aligned}[t]
  &\operatorname{Metrik}_{X}(d)\dsep\operatorname{Vollst}(X,d)
    \dsep A\subseteq X\dsep\operatorname{Abgeschlossen}_{d}(A)
    \vdash\\[-2pt]
  &\quad\operatorname{Vollst}(A,d_A).
  \end{aligned}%
}{ClosedSubspaceOfCompleteMetricSpace}{%
  \DeltaRow{Mengen}{X\dsep A}
  \DeltaRow{Funktionen}{d\dsep d_A}
}
\begin{proof}
Nach \FormulaRefAuto{MetricSubspaceRestrictionIsMetric}[thm] ist \(d_A\) eine
Metrik auf \(A\). Sei \(a\colon\mathbb N\to A\) eine Cauchy-Folge bezüglich
\(d_A\). Nach \FormulaRefAuto{MetricSubspaceCauchyAgree}[thm] ist sie als
Folge in \(X\) ebenfalls Cauchy und besitzt wegen der Vollständigkeit einen
Grenzwert \(x\in X\). Aus
\FormulaRefAuto{MetricClosedSetContainsSequenceLimits}[thm] folgt \(x\in A\).
Schließlich zeigt \FormulaRefAuto{MetricSubspaceConvergenceAgree}[thm], dass
die Folge auch bezüglich der Teilraummetrik \(d_A\) gegen \(x\) konvergiert.
\end{proof}

\chapter{Stetige Abbildungen}

\section{Stetigkeit in einem Punkt}

\FormulaDefDeltaK[Metrische Stetigkeit]{%
  \begin{aligned}[t]
  &\operatorname{Metrik}_{X}(d_X)\dsep
    \operatorname{Metrik}_{Y}(d_Y)\dsep
    f\colon X\to Y\dsep x\in X\vdash\\[-2pt]
  &\quad\operatorname{Stetig}_{d_X,d_Y}(f,x)\coloneqq
    \forall\varepsilon\in\mathbb R\,
    \Bigl(0<\varepsilon\rightarrow
      \exists\delta\in\mathbb R\,\bigl(0<\delta\land{}\\[-2pt]
  &\hspace{34mm}
      \forall u\in X\,
      (d_X(u,x)<\delta\rightarrow
       d_Y(f(u),f(x))<\varepsilon)\bigr)\Bigr),\\[2pt]
  &\operatorname{Metrik}_{X}(d_X)\dsep
    \operatorname{Metrik}_{Y}(d_Y)\dsep
    f\colon X\to Y\vdash\\[-2pt]
  &\quad\operatorname{Stetig}_{d_X,d_Y}(f)\coloneqq
    \forall x\in X\,\operatorname{Stetig}_{d_X,d_Y}(f,x).
  \end{aligned}%
}{MetricContinuityDef}{%
  \DeltaRow{Mengen}{X\dsep Y\dsep x\dsep u\dsep\varepsilon\dsep\delta}
  \DeltaRow{Funktionen}{d_X\dsep d_Y\dsep f}
  \DeltaRow{Neues Prädikat}{\operatorname{Stetig}_{d_X,d_Y}}
}

Das Bild eines Punktes soll also nahe bei \(f(x)\) liegen, sobald der Punkt
selbst hinreichend nahe bei \(x\) liegt. Der zulässige Eingangsradius
\(\delta\) darf von \(x\) und vom gewünschten Ausgangsradius
\(\varepsilon\) abhängen.

\FormulaThmDeltaK[Identität und konstante Abbildungen sind stetig]{%
  \begin{aligned}[t]
  &\operatorname{Metrik}_{X}(d_X)\dsep
    \operatorname{Metrik}_{Y}(d_Y)\dsep y_0\in Y\vdash\\[-2pt]
  &\quad\operatorname{Stetig}_{d_X,d_X}(\Id_X)\land
    \operatorname{Stetig}_{d_X,d_Y}
      (\operatorname{const}_{X,y_0}).
  \end{aligned}%
}{MetricIdentityConstantContinuous}{%
  \DeltaRow{Mengen}{X\dsep Y\dsep y_0}
  \DeltaRow{Funktionen}{d_X\dsep d_Y\dsep\Id_X\dsep
    \operatorname{const}_{X,y_0}}
}
\begin{proof}
Für die Identität kann zu \(\varepsilon>0\) derselbe Radius
\(\delta=\varepsilon\) verwendet werden. Bei der konstanten Abbildung ist
der Abstand der beiden Bildpunkte stets \(d_Y(y_0,y_0)=0\); jeder positive
Eingangsradius ist geeignet.
\end{proof}

\FormulaThmDeltaK[Der Abstand von einem festen Punkt ist stetig]{%
  \begin{aligned}[t]
  &\operatorname{Metrik}_{X}(d)\dsep z\in X\dsep
    \rho_z\colon X\to\mathbb R\dsep
    \forall u\in X\,\bigl(\rho_z(u)=d(u,z)\bigr)\vdash\\[-2pt]
  &\quad\operatorname{Stetig}_{d,d_{\mathbb R}}(\rho_z).
  \end{aligned}%
}{MetricDistanceFromPointContinuous}{%
  \DeltaRow{Mengen}{X\dsep z\dsep x\dsep u\dsep\varepsilon}
  \DeltaRow{Funktionen}{d\dsep d_{\mathbb R}\dsep\rho_z}
}
\begin{proof}
Fixiere \(x\in X\) und übernimm zu \(\varepsilon>0\) den Radius
\(\delta=\varepsilon\). Aus \(d(u,x)<\delta\) folgt mit der umgekehrten
Dreiecksungleichung
\[
 d_{\mathbb R}\bigl(\rho_z(u),\rho_z(x)\bigr)
 =|d(u,z)-d(x,z)|
 \leq d(u,x)<\varepsilon.
\]
Damit ist \(\rho_z\) in jedem Punkt \(x\) stetig.
\end{proof}

\FormulaThmDeltaK[Stetige Abbildungen erhalten Folgengrenzwerte]{%
  \begin{aligned}[t]
  &\operatorname{Metrik}_{X}(d_X)\dsep
    \operatorname{Metrik}_{Y}(d_Y)\dsep
    f\colon X\to Y\dsep x\in X\dsep
    \operatorname{Stetig}_{d_X,d_Y}(f,x)\dsep{}\\[-2pt]
  &\quad a\colon\mathbb N\to X\dsep a_n\MetricTo{d_X}x
    \vdash
    f(a_n)\MetricTo{d_Y}f(x).
  \end{aligned}%
}{MetricContinuousMapsPreserveLimits}{%
  \DeltaRow{Mengen}{X\dsep Y\dsep x\dsep\varepsilon\dsep\delta\dsep N}
  \DeltaRow{Funktionen}{d_X\dsep d_Y\dsep f\dsep a}
}
\begin{proof}
Zu \(\varepsilon>0\) liefert die Stetigkeit einen Radius \(\delta>0\).
Die Konvergenz von \(a\) liefert anschließend einen Index \(N\), ab dem
\(d_X(a_n,x)<\delta\) gilt. Dann ist für alle \(n\geq N\) auch
\(d_Y(f(a_n),f(x))<\varepsilon\).
\end{proof}

\section{Das Folgenkriterium}

Für die Rückrichtung verwenden wir die bereits bewiesene Nullfolge
\(1/(n+1)\), um eine Folge von Gegenbeispielen in immer kleineren Kugeln zu
konstruieren.

\FormulaThmDeltaK[Folgenkriterium für metrische Stetigkeit]{%
  \begin{aligned}[t]
  &\operatorname{Metrik}_{X}(d_X)\dsep
    \operatorname{Metrik}_{Y}(d_Y)\dsep
    f\colon X\to Y\dsep x\in X\vdash\\[-2pt]
  &\quad\operatorname{Stetig}_{d_X,d_Y}(f,x)
    \leftrightarrow
    \forall a\colon\mathbb N\to X\,
      \bigl(a_n\MetricTo{d_X}x
      \rightarrow f(a_n)\MetricTo{d_Y}f(x)\bigr).
  \end{aligned}%
}{MetricSequentialContinuityCriterion}{%
  \DeltaRow{Mengen}{X\dsep Y\dsep x\dsep\varepsilon\dsep n}
  \DeltaRow{Funktionen}{d_X\dsep d_Y\dsep f\dsep a\dsep r}
}
\begin{proof}
Die Hinrichtung ist
\FormulaRefAuto{MetricContinuousMapsPreserveLimits}[thm].

Für die Rückrichtung nehme man an, \(f\) sei in \(x\) nicht stetig. Dann
gibt es ein \(\varepsilon_0>0\), so dass für jedes \(\delta>0\) ein
\(u\in X\) mit
\[
 d_X(u,x)<\delta,
 \qquad
 d_Y(f(u),f(x))\geq\varepsilon_0
\]
existiert. Definiere die reelle Folge \(r\colon\mathbb N\to\mathbb R\) durch
\(r_n=1/(n+1)\). Nach
\FormulaRefAuto{ReciprocalNaturalSequenceConvergesToZero}[thm] gilt
\(r_n\to0\). Setzt man \(\delta=r_n\), so ist die entsprechende Menge
zulässiger Punkte für jedes \(n\) nichtleer. Das Auswahlprinzip
\FormulaRefAuto{Auswahlaxiom}[ax] liefert eine Folge \(a\colon\mathbb N\to X\)
mit
\[
 d_X(a_n,x)<r_n,
 \qquad
 d_Y(f(a_n),f(x))\geq\varepsilon_0.
\]
Zu jedem \(\varepsilon>0\) gilt ab einem Index \(N\) die Ungleichung
\(r_n<\varepsilon\); dann ist auch
\(d_X(a_n,x)<r_n<\varepsilon\). Somit konvergiert \(a\) gegen \(x\), aber die
Bildfolge konvergiert nicht gegen \(f(x)\). Dies widerspricht der rechten
Seite.
\end{proof}

\FormulaThmDeltaK[Komposition stetiger Abbildungen]{%
  \begin{aligned}[t]
  &\operatorname{Metrik}_{X}(d_X)\dsep
    \operatorname{Metrik}_{Y}(d_Y)\dsep
    \operatorname{Metrik}_{Z}(d_Z)\dsep
    f\colon X\to Y\dsep g\colon Y\to Z\dsep{}\\[-2pt]
  &\quad\operatorname{Stetig}_{d_X,d_Y}(f)\dsep
    \operatorname{Stetig}_{d_Y,d_Z}(g)
    \vdash
    \operatorname{Stetig}_{d_X,d_Z}(g\circ f).
\end{aligned}%
}{MetricContinuousComposition}{%
  \DeltaRow{Mengen}{X\dsep Y\dsep Z\dsep x\dsep u\dsep v\dsep
    \varepsilon\dsep\eta\dsep\delta}
  \DeltaRow{Funktionen}{d_X\dsep d_Y\dsep d_Z\dsep f\dsep g\dsep g\circ f}
}
\begin{proof}
Fixiere \(x\in X\) und \(\varepsilon>0\). Aus der Stetigkeit von \(g\) in
\(f(x)\) erhält man einen Radius \(\eta>0\), so dass
\[
 \forall v\in Y\,
 \Bigl(d_Y(v,f(x))<\eta
 \longrightarrow
 d_Z(g(v),g(f(x)))<\varepsilon\Bigr).
\]
Die Stetigkeit von \(f\) in \(x\) liefert zu diesem \(\eta\) einen Radius
\(\delta>0\), so dass
\[
 \forall u\in X\,
 \Bigl(d_X(u,x)<\delta
 \longrightarrow
 d_Y(f(u),f(x))<\eta\Bigr).
\]
Zusammen folgt für jedes \(u\in X\) mit \(d_X(u,x)<\delta\) die Abschätzung
\(d_Z((g\circ f)(u),(g\circ f)(x))<\varepsilon\). Dies ist die
Stetigkeit von \(g\circ f\) in \(x\); da \(x\) beliebig war, gilt sie auf
ganz \(X\).
\end{proof}

\section{Isometrien}

\FormulaDefDeltaK[Isometrische Abbildung]{%
  \begin{aligned}[t]
  &\operatorname{Metrik}_{X}(d_X)\dsep
    \operatorname{Metrik}_{Y}(d_Y)\dsep f\colon X\to Y\vdash\\[-2pt]
  &\quad\operatorname{Isometrie}_{d_X,d_Y}(f)\coloneqq
    \forall x,u\in X\,
      \bigl(d_Y(f(x),f(u))=d_X(x,u)\bigr).
  \end{aligned}%
}{MetricIsometryDef}{%
  \DeltaRow{Mengen}{X\dsep Y\dsep x\dsep u}
  \DeltaRow{Funktionen}{d_X\dsep d_Y\dsep f}
  \DeltaRow{Neues Prädikat}{\operatorname{Isometrie}_{d_X,d_Y}}
}

\FormulaThmDeltaK[Isometrien erhalten die metrische Struktur]{%
  \begin{aligned}[t]
  &\operatorname{Metrik}_{X}(d_X)\dsep
    \operatorname{Metrik}_{Y}(d_Y)\dsep f\colon X\to Y\dsep{}\\[-2pt]
  &\quad\operatorname{Isometrie}_{d_X,d_Y}(f)\vdash
    \Bigl(\operatorname{Stetig}_{d_X,d_Y}(f)\land{}\\[-2pt]
  &\qquad\forall a\colon\mathbb N\to X\,
    \bigl(\operatorname{Cauchy}_{d_X}(a)
      \rightarrow\operatorname{Cauchy}_{d_Y}(f\circ a)\bigr)\Bigr).
  \end{aligned}%
}{MetricIsometryPreservesStructure}{%
  \DeltaRow{Mengen}{X\dsep Y}
  \DeltaRow{Funktionen}{d_X\dsep d_Y\dsep f\dsep a\dsep f\circ a}
}
\begin{proof}
In beiden Definitionen werden nur Abstände verglichen. Eine Isometrie
ersetzt jeden Abstand im Urbildraum durch denselben reellen Wert im
Bildraum. Für die Stetigkeit kann daher \(\delta=\varepsilon\) gewählt
werden; dieselbe Gleichheit überträgt die Cauchy-Ungleichungen.
\end{proof}

\chapter{Folgenkompaktheit}

\section{Definition und erste Folgerungen}

\FormulaDefDeltaK[Folgenkompakte Teilmenge]{%
  \begin{aligned}[t]
  &\operatorname{Metrik}_{X}(d)\dsep K\subseteq X\vdash
    \operatorname{Folgenkompakt}_{d}(K)\coloneqq{}\\[-2pt]
  &\quad\forall a\colon\mathbb N\to K\,
    \exists x\in K\,\exists\varphi\colon\mathbb N\to\mathbb N\,
    \Bigl(\forall n\in\mathbb N\,
      \bigl(\varphi(n)<\varphi(\Succ(n))\bigr)\\[-2pt]
  &\hspace{64mm}\land
      a_{\varphi(n)}\MetricTo{d}x\Bigr).
  \end{aligned}%
}{MetricSequentialCompactnessDef}{%
  \DeltaRow{Mengen}{X\dsep K\dsep x\dsep n}
  \DeltaRow{Funktionen}{d\dsep a\dsep\varphi}
  \DeltaRow{Neues Prädikat}{\operatorname{Folgenkompakt}_{d}}
}

Der Grenzwert muss in \(K\) liegen. Eine bloß in \(X\) konvergente Teilfolge
genügt daher nicht. Der Begriff wird zunächst bewusst als
\emph{Folgenkompaktheit} bezeichnet; die spätere Gleichwertigkeit mit der
Überdeckungskompaktheit ist ein eigener topologischer Satz.

\FormulaThmDeltaK[Wachsende Aufzählung unendlicher Mengen natürlicher Zahlen]{%
  \begin{aligned}[t]
  &J\subseteq\mathbb N\dsep\Infinite{J}\vdash
    \exists\varphi\colon\mathbb N\to\mathbb N\,\forall n\in\mathbb N\,{}\\[-2pt]
  &\quad\Bigl(\varphi(n)\in J\land
    \varphi(n)<\varphi(\Succ(n))\Bigr).
  \end{aligned}%
}{InfiniteNaturalSubsetIncreasingEnumeration}{%
  \DeltaRow{Mengen}{J\dsep p\dsep n}
  \DeltaRow{Funktionen}{s_J\dsep\varphi}
}
\begin{proof}
Die Menge \(J\) ist nichtleer. Außerdem ist sie nach oben unbeschränkt:
Wären alle ihre Elemente höchstens \(p\in\mathbb N\), so wäre
\(J\subseteq\NatSeg{\Succ(p)}\) und damit nach
\FormulaRefAuto{FiniteNatSegSubset}[thm] endlich, ein Widerspruch.

Somit ist für jedes \(p\in\mathbb N\) die Menge
\[
 J_{>p}\coloneqq\{j\in J\mid p<j\}
\]
nichtleer. Das Wohlordnungsprinzip
\FormulaRefAuto{NaturalWellOrderingPrinciple}[thm] liefert ihr jeweils ein
eindeutiges kleinstes Element. Das Prinzip der typisierten
Funktionsdefinition liefert daher eine Funktion
\(s_J\colon\mathbb N\to\mathbb N\) mit
\(s_J(p)=\min J_{>p}\).

Wähle ferner \(j_0=\min J\). Der Dedekindsche Rekursionssatz
\FormulaRefAuto{DedekindRecursionTheorem}[thm] liefert eine Funktion
\(\varphi\colon\mathbb N\to\mathbb N\) mit
\[
 \varphi(0)=j_0,
 \qquad
 \varphi(\Succ(n))=s_J(\varphi(n)).
\]
Nach Konstruktion liegen alle Werte in \(J\), und jeder Nachfolgerwert ist
strikt größer als sein Vorgänger.
\end{proof}

\FormulaThmDeltaK[Endliche Teilmengen sind folgenkompakt]{%
  \operatorname{Metrik}_{X}(d)\dsep K\subseteq X\dsep\Finite{K}
  \vdash\operatorname{Folgenkompakt}_{d}(K)%
}{FiniteMetricSetSequentiallyCompact}{%
  \DeltaRow{Mengen}{X\dsep K\dsep x\dsep n}
  \DeltaRow{Funktionen}{d\dsep a\dsep\varphi}
}
\begin{proof}
Sei \(a\colon\mathbb N\to K\). Mindestens ein Wert \(x\in K\) wird
unendlich oft angenommen. Denn wäre jede Faser \(a^{-1}[\{x\}]\) endlich,
so zeigte die endliche Induktion über \(K\) gemäß
\FormulaRefAuto{FiniteInductionRule}[ax], dass auch ihre Vereinigung
endlich ist: Im Induktionsanfang ist sie leer, und beim Hinzufügen eines
Punktes kommt nach \FormulaRefAuto{FiniteUnion}[thm] nur eine weitere endliche
Faser hinzu. Wegen
\[
 \mathbb N=\bigcup_{x\in K}a^{-1}[\{x\}]
\]
widerspräche dies
\FormulaRefAuto{NaturalNumbersInfinite}[thm]. Also ist für ein \(x\in K\) die
Indexmenge
\[
 J\coloneqq\{n\in\mathbb N\mid a_n=x\}
\]
unendlich. Nach
\FormulaRefAuto{InfiniteNaturalSubsetIncreasingEnumeration}[thm] besitzt sie
eine streng wachsende Aufzählung \(\varphi\). Damit gilt
\(a_{\varphi(n)}=x\) für alle \(n\), und die konstante Teilfolge konvergiert
gegen \(x\in K\).
\end{proof}

\FormulaThmDeltaK[Folgenkompakte Teilmengen sind beschränkt]{%
  \operatorname{Metrik}_{X}(d)\dsep K\subseteq X\dsep
  \operatorname{Folgenkompakt}_{d}(K)
  \vdash\operatorname{Beschr}_{d}(K)%
}{MetricSequentiallyCompactSetBounded}{%
  \DeltaRow{Mengen}{X\dsep K\dsep c\dsep n}
  \DeltaRow{Funktionen}{d\dsep a\dsep\varphi}
}
\begin{proof}
Für \(K=\varnothing\) gilt die Behauptung nach Definition. Sei also
\(c\in K\). Wäre \(K\) unbeschränkt, so gäbe es für jedes
  \(n\in\mathbb N\) einen Punkt \(a_n\in K\) mit
\(d(a_n,c)\geq n+1\). Das Auswahlprinzip
\FormulaRefAuto{Auswahlaxiom}[ax] liefert daraus eine Folge
\(a\colon\mathbb N\to K\). Dies ist die einzige Auswahlstelle dieses
Beweises.

Nach Folgenkompaktheit besitzt sie eine gegen einen Punkt von \(K\)
konvergente Teilfolge \(a_{\varphi(n)}\). Diese Teilfolge ist nach
\FormulaRefAuto{ConvergentMetricSequenceBounded}[thm] beschränkt. Andererseits
gilt wegen \(\varphi(n)\geq n\)
\[
 d(a_{\varphi(n)},c)\geq\varphi(n)+1\geq n+1,
\]
Die archimedische Eigenschaft
\FormulaRefAuto{RealArchimedeanProperty}[thm] liefert zu jedem reellen Radius
\(r>0\) ein \(n\) mit \(r<n+1\). Daher liegt die Teilfolge in keiner Kugel um
\(c\). Jede Kugel um einen anderen Mittelpunkt liegt nach der
Dreiecksungleichung in einer hinreichend großen Kugel um \(c\). Damit ist die
Teilfolge unbeschränkt, ein Widerspruch.
\end{proof}

\FormulaThmDeltaK[Folgenkompakte Teilräume sind vollständig]{%
  \operatorname{Metrik}_{X}(d)\dsep K\subseteq X\dsep
  \operatorname{Folgenkompakt}_{d}(K)
  \vdash\operatorname{Vollst}(K,d_K)%
}{MetricSequentiallyCompactSetComplete}{%
  \DeltaRow{Mengen}{X\dsep K}
  \DeltaRow{Funktionen}{d\dsep d_K\dsep a\dsep\varphi}
}
\begin{proof}
Nach \FormulaRefAuto{MetricSubspaceRestrictionIsMetric}[thm] ist \(d_K\) eine
Metrik auf \(K\). Sei \(a\colon\mathbb N\to K\) eine Cauchy-Folge bezüglich
\(d_K\). Nach \FormulaRefAuto{MetricSubspaceCauchyAgree}[thm] ist sie auch
bezüglich \(d\) Cauchy. Die Folgenkompaktheit liefert eine bezüglich \(d\)
gegen ein \(x\in K\) konvergente Teilfolge. Nach
\FormulaRefAuto{MetricSubspaceConvergenceAgree}[thm] konvergiert diese
Teilfolge auch bezüglich \(d_K\). Nun ist
\FormulaRefAuto{MetricCauchyWithConvergentSubsequence}[thm] im metrischen Raum
\((K,d_K)\) anwendbar und zeigt, dass die ganze Folge gegen \(x\) konvergiert.
\end{proof}

\FormulaThmDeltaK[Folgenkompakte Teilmengen sind abgeschlossen]{%
  \operatorname{Metrik}_{X}(d)\dsep K\subseteq X\dsep
  \operatorname{Folgenkompakt}_{d}(K)
  \vdash\operatorname{Abgeschlossen}_{d}(K)%
}{MetricSequentiallyCompactSetClosed}{%
  \DeltaRow{Mengen}{X\dsep K\dsep x\dsep y\dsep n}
  \DeltaRow{Funktionen}{d\dsep a\dsep r\dsep\varphi}
}
\begin{proof}
Angenommen, \(K\) sei nicht abgeschlossen. Dann ist \(X\setminus K\) nicht
offen. Folglich gibt es ein \(x\in X\setminus K\), so dass jede positive
Kugel um \(x\) einen Punkt aus \(K\) enthält. Definiere die reelle Folge
\(r\colon\mathbb N\to\mathbb R\) durch \(r_n=1/(n+1)\). Dann ist für jedes
\(n\in\mathbb N\) die Menge
\[
 K\cap\MetricBall{d}{x}{r_n}
\]
nichtleer. Das Auswahlprinzip
\FormulaRefAuto{Auswahlaxiom}[ax] liefert eine Folge
\(a\colon\mathbb N\to K\) mit
\(d(a_n,x)<r_n\). Nach
\FormulaRefAuto{ReciprocalNaturalSequenceConvergesToZero}[thm] gilt
\(r_n\to0\). Zu \(\varepsilon>0\) ist daher ab einem Index \(N\)
\(r_n<\varepsilon\), also auch \(d(a_n,x)<\varepsilon\). Somit gilt
\(a_n\MetricTo{d}x\).

Die Folgenkompaktheit liefert eine Teilfolge \(a_{\varphi(n)}\), die gegen
ein \(y\in K\) konvergiert. Nach
\FormulaRefAuto{MetricSubsequencePreservesLimit}[thm] konvergiert dieselbe
Teilfolge auch gegen \(x\). Die Eindeutigkeit metrischer Grenzwerte ergibt
\(x=y\in K\), im Widerspruch zu \(x\notin K\).
\end{proof}

\FormulaThmDeltaK[Stetige Bilder folgenkompakter Mengen]{%
  \begin{aligned}[t]
  &\operatorname{Metrik}_{X}(d_X)\dsep
    \operatorname{Metrik}_{Y}(d_Y)\dsep f\colon X\to Y\dsep
    K\subseteq X\dsep{}\\[-2pt]
  &\quad\operatorname{Stetig}_{d_X,d_Y}(f)\dsep
    \operatorname{Folgenkompakt}_{d_X}(K)
    \vdash
    \operatorname{Folgenkompakt}_{d_Y}(f[K]).
  \end{aligned}%
}{MetricContinuousImageSequentiallyCompact}{%
  \DeltaRow{Mengen}{X\dsep Y\dsep K}
  \DeltaRow{Funktionen}{d_X\dsep d_Y\dsep f\dsep a\dsep b\dsep\varphi}
}
\begin{proof}
Sei \(b\colon\mathbb N\to f[K]\). Für jedes \(n\) ist die Menge
\(K\cap f^{-1}[\{b_n\}]\) nichtleer. Das Auswahlprinzip liefert eine Folge
\(a\colon\mathbb N\to K\) mit \(f(a_n)=b_n\); formal wird hier
\FormulaRefAuto{Auswahlaxiom}[ax] verwendet. Aus der Folgenkompaktheit von
\(K\) folgt eine Teilfolge \(a_{\varphi(n)}\), die gegen ein \(x\in K\)
konvergiert. Die Stetigkeit und
\FormulaRefAuto{MetricContinuousMapsPreserveLimits}[thm] liefern
\[
 b_{\varphi(n)}=f(a_{\varphi(n)})\MetricTo{d_Y}f(x),
\]
wobei \(f(x)\in f[K]\) gilt.
\end{proof}

\section{Warum Beschränktheit allein nicht genügt}

\FormulaThmDeltaK[Konvergenz in der diskreten Metrik]{%
  a\colon\mathbb N\to X\dsep x\in X
  \vdash
  \bigl(a_n\MetricTo{\delta_X}x\bigr)
  \leftrightarrow
  \exists N\in\mathbb N\,\forall n\in\mathbb N\,
    (N\leq n\rightarrow a_n=x)%
}{DiscreteMetricConvergenceEventuallyConstant}{%
  \DeltaRow{Mengen}{X\dsep x\dsep N\dsep n}
  \DeltaRow{Funktionen}{a\dsep\delta_X}
}
\begin{proof}
Konvergiert \(a\) gegen \(x\), so gilt ab einem Index
\(\delta_X(a_n,x)<1/2\). Die diskrete Metrik kann dann nicht den Wert \(1\)
haben, also ist \(a_n=x\). Die Rückrichtung folgt unmittelbar, weil eine
schließlich konstante Folge gegen ihren konstanten Wert konvergiert.
\end{proof}

\FormulaThmDeltaK[Eine beschränkte, nicht folgenkompakte Metrik]{%
  \operatorname{Beschr}_{\delta_{\mathbb N}}(\mathbb N)
  \land
  \neg\operatorname{Folgenkompakt}_{\delta_{\mathbb N}}(\mathbb N)%
}{BoundedDiscreteNaturalNumbersNotSequentiallyCompact}{%
  \DeltaRow{Mengen}{\mathbb N}
  \DeltaRow{Funktionen}{\delta_{\mathbb N}}
}
\begin{proof}
Die Kugel \(B_2^{\delta_{\mathbb N}}(0)\) ist ganz \(\mathbb N\), also ist
der Raum beschränkt. Die kanonische Folge \(a_n=n\) besitzt jedoch keine
konvergente Teilfolge. Jede Teilfolge mit streng wachsenden Indizes nimmt
weiterhin paarweise verschiedene Werte an und ist daher nicht schließlich
konstant. Nach
\FormulaRefAuto{DiscreteMetricConvergenceEventuallyConstant}[thm] kann sie
nicht konvergieren.
\end{proof}

Dieses Beispiel markiert eine entscheidende Grenze der Verallgemeinerung:
In einem beliebigen metrischen Raum folgt aus Beschränktheit keine
konvergente Teilfolge. Der Satz von Bolzano--Weierstraß ist daher kein bloßer
Metriksatz.

\section{Bolzano--Weierstraß auf den reellen Zahlen}

\FormulaThmDeltaK[Satz von der monotonen Teilfolge]{%
  \begin{aligned}[t]
  &a\colon\mathbb N\to\mathbb R\vdash
    \exists\varphi\colon\mathbb N\to\mathbb N\,\Bigl({}\\[-2pt]
  &\quad\forall n\in\mathbb N\,
      \bigl(\varphi(n)<\varphi(\Succ(n))\bigr)\land{}\\[-2pt]
  &\quad\bigl(\operatorname{Mon}_{\uparrow}(a\circ\varphi)
    \lor
    \operatorname{Mon}_{\downarrow}(a\circ\varphi)\bigr)\Bigr).
  \end{aligned}%
}{RealSequenceMonotoneSubsequence}{%
  \DeltaRow{Mengen}{n\dsep m}
  \DeltaRow{Funktionen}{a\dsep\varphi\dsep a\circ\varphi}
}
\begin{proof}
Nenne einen Index \(n\) einen Spitzenindex, wenn
\[
  \forall m\in\mathbb N\,(n<m\rightarrow a_m\leq a_n).
\]
Gibt es unendlich viele Spitzenindizes, so liefert
\FormulaRefAuto{InfiniteNaturalSubsetIncreasingEnumeration}[thm] eine streng
wachsende Aufzählung \(\varphi\) dieser Indizes. Aus der Spitzeneigenschaft
folgt
\(a_{\varphi(\Succ(n))}\leq a_{\varphi(n)}\). Die Teilfolge ist monoton
fallend.

Gibt es nur endlich viele Spitzenindizes, so wähle man einen Index
\(N\), oberhalb dessen keiner mehr liegt. Jeder Index \(n\geq N\) besitzt
dann einen späteren Index \(m>n\) mit \(a_n<a_m\). Beginnend mit
\(\varphi(0)=N\) wählt man rekursiv den jeweils kleinsten solchen Index.
Dann ist sowohl \(\varphi\) als auch die Wertfolge
\(a_{\varphi(n)}\) streng wachsend.
\end{proof}

\FormulaThmDeltaK[Bolzano--Weierstraß für reelle Folgen]{%
  \begin{aligned}[t]
  &a\colon\mathbb N\to\mathbb R\dsep\operatorname{Beschr}(a)
    \vdash\\[-2pt]
  &\quad\exists L\in\mathbb R\,
    \exists\varphi\colon\mathbb N\to\mathbb N\,
    \Bigl(\forall n\in\mathbb N\,
      \bigl(\varphi(n)<\varphi(\Succ(n))\bigr)\\[-2pt]
  &\hspace{65mm}\land
      a_{\varphi(n)}\longrightarrow L\Bigr).
  \end{aligned}%
}{RealSequenceBolzanoWeierstrass}{%
  \DeltaRow{Mengen}{L\dsep n}
  \DeltaRow{Funktionen}{a\dsep\varphi\dsep a\circ\varphi}
}
\begin{proof}
Nach \FormulaRefAuto{RealSequenceMonotoneSubsequence}[thm] besitzt \(a\) eine
monotone Teilfolge. Die Beschränktheit von \(a\) vererbt sich an diese
Teilfolge. Ist sie wachsend, so ist sie nach oben beschränkt; ist sie
fallend, so ist sie nach unten beschränkt. In beiden Fällen liefert
\FormulaRefAuto{MonotoneRealSequenceConvergence}[thm] einen reellen Grenzwert.
\end{proof}

\FormulaThmDeltaK[Abgeschlossene beschränkte reelle Intervalle sind folgenkompakt]{%
  A,B\in\mathbb R\dsep A\leq B
  \vdash
  \operatorname{Folgenkompakt}_{d_{\mathbb R}}
    \bigl([A,B]_{\mathbb R}\bigr)%
}{RealClosedBoundedIntervalSequentiallyCompact}{%
  \DeltaRow{Mengen}{A\dsep B\dsep L\dsep n}
  \DeltaRow{Funktionen}{d_{\mathbb R}\dsep a\dsep\varphi}
}
\begin{proof}
Jede Folge in \([A,B]_{\mathbb R}\) ist reell beschränkt. Nach
\FormulaRefAuto{RealSequenceBolzanoWeierstrass}[thm] besitzt sie eine gegen
ein \(L\in\mathbb R\) konvergente Teilfolge. Für jedes Glied dieser
Teilfolge gilt \(A\leq a_{\varphi(n)}\leq B\). Die konstanten Folgen mit den
Werten \(A\) und \(B\) konvergieren gegen \(A\) beziehungsweise \(B\); daher
liefert
\FormulaRefAuto{RealSequenceLimitPreservesOrder}[thm]
\[
 A\leq L\leq B.
\]
Somit liegt der Grenzwert in \([A,B]_{\mathbb R}\). Nach
\FormulaRefAuto{RealAndMetricSequenceConvergenceAgree}[thm] ist dies zugleich
metrische Konvergenz.
\end{proof}

Der Vergleich fasst die Lage zusammen: In jedem metrischen Raum folgt aus
Folgenkompaktheit die Beschränktheit, aber nicht umgekehrt. Im besonderen
Raum \((\mathbb R,d_{\mathbb R})\) besitzt dagegen jede beschränkte Folge
eine konvergente Teilfolge; deshalb sind die abgeschlossenen beschränkten
Intervalle folgenkompakt.

\chapter{Ausblick: Normen}

Eine Norm misst zunächst die Größe eines einzelnen Vektors. Erst die
Vektorraumoperationen erlauben es, daraus den Abstand zweier Punkte durch die
Größe ihrer Differenz zu gewinnen. Ist \(V\) ein reeller Vektorraum, so
fordert eine Norm in der späteren formalen Definition
\[
\begin{aligned}
 &\|v\|\geq0,
 &&\|v\|=0\leftrightarrow v=0,\\
 &\|\lambda v\|=|\lambda|\,\|v\|,
 &&\|u+v\|\leq\|u\|+\|v\|.
\end{aligned}
\]
Dann erzeugt
\[
 d_{\|\cdot\|}(u,v)\coloneqq\|u-v\|
\]
eine Metrik: Positivität und Definitheit kommen aus den ersten beiden
Normgesetzen, Symmetrie aus
\(\|u-v\|=\|-(v-u)\|=\|v-u\|\), und die Dreiecksungleichung aus
\[
 \|u-w\|=\|(u-v)+(v-w)\|
 \leq\|u-v\|+\|v-w\|.
\]

Damit gilt die begriffliche Richtung
\[
 \text{Normraum}\quad\Longrightarrow\quad\text{metrischer Raum}.
\]
Die Umkehrung gilt nicht: Eine allgemeine Metrik kennt weder Vektoraddition
noch Skalarmultiplikation. Eine formale Normdefinition an dieser Stelle würde
daher Begriffe voraussetzen, die im bisherigen Aufbau noch nicht vorhanden
sind. Sie wird nach der Konstruktion der Vektorräume eingeführt; alle bis
dahin bewiesenen metrischen Sätze stehen dort sofort zur Verfügung.

Die nächsten Schritte der Analysis können nun auf einer klaren Grundlage
aufbauen: Reihen verwenden zunächst die Vollständigkeit von
\(\mathbb R\), Stetigkeit reeller Funktionen ist ein Spezialfall der hier
definierten metrischen Stetigkeit, und Kompaktheit kann ausgehend von der
Folgenkompaktheit abgeschlossener Intervalle entwickelt werden.

\end{document}
